\section*{Contents}
\begin{tabularx}{0.98\linewidth}{@{}>{\bfseries}p{0.28\linewidth}X@{}}
  1 & Frontier Models --- GPT-6とHaiku 5.5の提供条件 \\
  2 & Agent Infrastructure --- 検証済みAPIと観測・告知の層別 \\
  3 & Open Weights and Audio --- 公開言明と取得実態の7件 \\
  4 & Image Model Card --- Nano Banana 2.1の仕様読み \\
  5 & Provenance and Watermarks --- textGrainとSynthIDの対置 \\
  6 & Safety and Science --- 制度・公開・模擬の境界線 \\
  7 & Community Movement --- 公開投稿の関心と限界 \\
  8 & Week in Review --- 提供条件と検証範囲の総括 \\
\end{tabularx}

\vspace{4mm}
\section*{This Week in AI}

\begin{tcolorbox}[enhanced,colback=SurveySoft,colframe=SurveyRule,boxrule=0.4pt,arc=1mm,left=3mm,right=3mm,top=2mm,bottom=2mm]
\textbf{提供条件が具体化した。}\quad 10月7日のGPT-6とIntelligent UI、10月7日のHaiku 5.5、10月6日のDecisions API betaと10月8日のUltrafast tier、10月3日のKolibriと10月7日のd1、10月6日のEmbeddingGemma 2、10月5日のBeamと10月6日のMistral Large 4、10月9日のClef-omniと10月5日のNova Sonic、10月6日のNano Bananaカードと10月5日のtextGrain、10月7日のSynthID拡張、10月8日のCyber Missionと10月6日の数学公開・Ironclad模擬が並んだ。
\end{tcolorbox}

\begin{tcolorbox}[enhanced,colback=SurveySoft,colframe=SurveyRule,boxrule=0.4pt,arc=1mm,left=3mm,right=3mm,top=2mm,bottom=2mm]
\textbf{数値と日付は出どころごとに受け止める。}\quad ベンチマークの値はベンダー計測によるもので、独立した再現はない。料金表は発表時点の公式ページであり、請求額の実測ではない。公開の言明と配布先の提示、実ファイルの取得、重みの公開完了は別の事柄として切り分ける。透かしの検出率と回答品質、検証の射程と効き目は混ぜない。未検証の事柄、査読前の報告、構想段階のものは、その旨を添えて読む。
\end{tcolorbox}

\begin{tcolorbox}[enhanced,colback=SurveySoft,colframe=SurveyRule,boxrule=0.4pt,arc=1mm,left=3mm,right=3mm,top=2mm,bottom=2mm]
\textbf{載せないものは理由とともに記す。}\quad 持ち越しの4件は新しい日付のある証拠がなく今号の章には載せない。締め切り後の速報は次号へ送り、通常の集計には含めない。公開投稿の反応は関心の方向を示す背景にとどめる。
\end{tcolorbox}

\vspace{2mm}
\begin{claimboundary}[この号の読み方]
通常本文の対象期間は2026-10-02 18:00から2026-10-09 18:00（いずれもAmerica/New\_York、終了時刻は含まない）。この範囲の出来事を通常の扱いとし、期間前の持ち越し、締め切り後の速報、予告された今後の運用は区別する。本号で参照する公開資料が示す範囲を超えて断定しない。ベンダーや報道が伝える数値はそれぞれの主体に帰属させ、条件の違う測定値を横断的な順位へ言い換えない。モデルの公開場所や一次発表で確認できない仕様や利用条件は事実として扱わない。
\end{claimboundary}

\vspace{2mm}
\noindent\textit{読み解く軸：} 提供条件の確定と検証範囲の確定、公開言明と実取得の切り分け、ベンダー計測の帰属、透かし検出と検証射程の分離を、一次資料の範囲と留保とともに束ねることである。
